% Generated draft: requires mathematical and source review.
\documentclass[11pt]{article}
\usepackage[margin=25mm]{geometry}
\usepackage{fontspec}
\setmainfont{Latin Modern Roman}
\usepackage{amsmath,amssymb,mathtools,unicode-math}
\setmathfont{Latin Modern Math}
\usepackage{enumitem,needspace,xurl,hyperref}
\setlength{\parindent}{0pt}
\setlength{\parskip}{5pt}
\setlength{\emergencystretch}{3em}
\newcommand{\N}{\mathbb{N}}
\newcommand{\Z}{\mathbb{Z}}
\newcommand{\Q}{\mathbb{Q}}
\newcommand{\R}{\mathbb{R}}
\newcommand{\C}{\mathbb{C}}
\newcommand{\F}{\mathbb{F}}
\newcommand{\A}{\mathbb{A}}
\newcommand{\PP}{\mathbb{P}}
\newcommand{\E}{\mathbb{E}}
\newcommand{\eps}{\varepsilon}
\newcommand{\dd}{\,\mathrm{d}}
\begin{document}
\section*{UnsolvedMath ID: 20003097}
\textbf{Source problem code:} catalogue entry
\subsection{Definitions and mathematical statement}
The source asks to define a geodesic flow for the asymmetric Thurston metric on Teichm\"uller space and study ergodicity and mixing, inspired by Eskin--Mirzakhani's work for the Teichm\"uller metric. Nonuniqueness of Thurston geodesics means a flow requires a specified state space and geodesic selection. Ergodicity and measure-theoretic mixing additionally require an invariant measure; marked space and its mapping-class quotient are different dynamical settings. The extracted following sentence and page header are not additional hypotheses of Problem 2.7.

The supplied aid reports Pan--Wolf's 2026 continuous endpoint-labelled harmonic-stretch selection. Write \(\mathcal P\) for its state space, \(\psi_t\) for the selected flow, \(X(p)\in\mathcal T(S)\) for the marked basepoint of a state, and \(\eta(p)\) for its forward endpoint, unchanged by the flow. The complete definition, surface range and domain of this state space are not supplied here; this notation is conditional on that construction, not a claim that every tangent vector canonically determines a geodesic.

Fix \(O\in\mathcal T(S)\). Use the forward Walsh horofunction normalization
\[
h_\eta^O(X)=\lim_{Z\to\eta}\bigl(d_{\rm Th}(X,Z)-d_{\rm Th}(O,Z)\bigr),\qquad h_\eta^O(O)=0,
\]
where the limit is in the relevant horofunction compactification. The aid reports that on the selected endpoint flow,
\[
B_O(p)=h_{\eta(p)}^O(X(p)),\qquad B_O(\psi_t p)=B_O(p)-t.
\]
Thus \(B_O\) is an exact continuous time coordinate on marked state space. No state is forward recurrent: returns to a fixed neighborhood with bounded \(B_O\)-values are impossible for sufficiently large positive time. The same coordinate excludes topological mixing and any nonzero finite invariant Borel measure. Indeed such a measure would push forward to a nonzero finite translation-invariant measure on \(\R\), which cannot exist. These statements concern the marked flow and do not address infinite invariant measures or assert quotient nonergodicity.

For a mapping class \(g\), the normalized horofunction transforms as
\[
B_O(gp)=B_O(p)-h_{\eta(p)}^O(g^{-1}O).
\]
The endpoint term is a basepoint-dependent cocycle, so \(B_O\) need not descend to the quotient. If quotient recurrence is witnessed by \(t_n\to+\infty\) and mapping classes \(g_n\) with \(g_n\psi_{t_n}p\to p\), the identity forces
\[
h_{\eta(p)}^O(g_n^{-1}O)=-t_n+o(1).
\]
Thus quotient recurrence requires cocycle compensation of the marked drift; the marked obstruction alone cannot prohibit it. General quotient ergodicity, mixing, and the intended invariant-measure construction remain open in the supplied account. The problem is not resolved simply by selecting geodesics or by showing nonrecurrence before taking the quotient.
\subsection{Short English statement}
Choose a precise Thurston geodesic flow and a measure before asking about ergodicity or mixing. The supplied endpoint-selected flow has a horofunction drifting exactly by minus time on marked space; a mapping-class cocycle prevents this obstruction from automatically descending to the quotient.
\subsection{Sources}
\begin{itemize}
\item K. Rafi, Problem 2.7, AIM \emph{Problems on the Thurston metric}, workshop \emph{Lipschitz metric on Teichmueller space}, \url{https://aimath.org/pastworkshops/lipschitzteichproblems.pdf}. The source cites Eskin--Mirzakhani's Teichm\"uller-metric work as motivation without a full reference.
\item The research summary cites Pan--Wolf's 2026 endpoint-labelled harmonic-stretch construction and Walsh horofunctions, without complete publication titles, URLs, theorem statements or state-space definitions.
\item Local UnsolvedMath registry, permanent ID 20003097, exact code \texttt{catalogue entry}; canonical location \texttt{aim-topology-notes.json}, notes[184], dated 14 August 2026; statement labeled exact, with trailing extracted transition/header text separated here.
\item Research attempt 1: partial result, problem partially solved at the run, rigor label proved, novelty confidence low; registry status partially solved. The drift and cocycle reductions are machine-generated research aids, not peer reviewed; independent verification is required before citation.
\end{itemize}
\end{document}
